\subsubsection{The archimedean jet of a single kernel}
\label{subsec:selected-kernel-jets}

For a single term in the class sum one must bound a derivative of the archimedean kernel, including both jets when the last two lifts are coupled.

\paragraph{The padded polynomial bound.}

The first estimate bounds the archimedean polynomial after extending it to the larger of the two ranks in a lift.

\begin{established}[Determinant shifts in the padded jet model]\label{hyp:padded-Gaussian-jets}
For integers \(1\le n\le n'\), let \(\nu_\tau\) be a partition of
\(q_\tau\) of length at most \(n\) at each real label \(\tau\).
On the source divide out the common last entry, and on the target
use the partition padded with zeros to length \(n'\), up to its
determinant character. With the highest-vector metrics, restoring
the source determinant factor by multiplication with
\(\det^{\nu_{\tau n}}\) gives an isometry \(T\). In particular
\[
               T^*T=I,\qquad \deg(Tv)=|\nu_\tau|=q_\tau .
\]
\end{established}
\begin{proof}
In the flag-minor model the determinant has weight \((1,\ldots,1)\)
and has absolute value one on the unitary group. Thus \(T^*T\)
commutes with that group. Schur's lemma makes it scalar, and the
normalized source highest polynomial is sent to the normalized
target highest polynomial. The scalar is their squared norm ratio,
which is one. If \(n=n'\), padding inserts no zeros. The common
last entry is removed by the determinant twist in this case as well.
\end{proof}

\begin{lemma}[Norm of a padded jet contraction]\label{lem:padded-jet-estimate}
Use the polynomial spaces of Lemma~\ref{hyp:padded-Gaussian-jets}.
For \(Z_\tau\in\mathbb C^{n'\times n}\) substitute
\(Y_\tau=Z_\tau^t X_\tau E_n\), where \(E_n:\mathbb C^n\hookrightarrow\mathbb C^{n'}\)
includes the first coordinates and \(X_\tau\) is the square target
variable. Let \(G\) be entire, with \(S\ge0\), \(r_\tau,\Omega_\tau>0\) and
\[
 |G(Z)|\le S\exp\!\left(\tfrac12\sum_\tau r_\tau^2\|Z_\tau\|^2\right),
 \qquad |r_\tau d_\tau|^2=\Omega_\tau^{-1}.
\]
Here the matrix norm is Euclidean on the entries.
The jet replaces \(Z^\alpha\) by
\((\prod_\tau d_\tau^{|\alpha_\tau|})\partial^\alpha G(0)\).
The squared operator norm of this contraction,
in the normalized highest-vector metrics, is at most
\begin{equation}\label{eq:selected-padded-constant}
 S^2\prod_\tau\left[
    \Omega_\tau^{-q_\tau} e(2q_\tau+2)^{2nn'}
    \exp\!\left(\sum_i\varphi(\nu_{\tau i})\right)
 \right].
\end{equation}
This bound includes labels of degree zero.
\end{lemma}

\begin{proof}
Write \(W_\nu(\C^d)\) for the Schur module with its unitary
metric, and \(v_{\nu,d}\) for its initial flag vector of norm
one. The realization \(u\mapsto u^t W_\nu(Y)v_{\nu,d}\)
is the flag-minor model, with the transpose taken in bases
compatible with the real structures. In fact the highest
matrix coefficient is the product of initial flag minors,
and the row action generates the model. The Schur maps
commute with transpose in these bases. Functoriality and compatibility with transpose give
\[
 u^tW_\nu(Z^t X E_n)v_{\nu,n}
 =(W_\nu(Z)u)^t W_\nu(X) W_\nu(E_n)v_{\nu,n}
 =(W_\nu(Z)u)^t W_\nu(X)v_{\nu,n'} .
\] The identity uses
\(W_\nu(E_n)v_{\nu,n}=v_{\nu,n'}\), with compatible highest
vectors. Thus the coefficient contraction lands in the
target highest-weight model. This explains the order
\(Z^tX E_n\) in the substitution and the absence of a
projection cost.

\proofstep{Gaussian contraction of the two representations.}
Use the Gaussian polynomial norm
\(\|p\|_{\rm F}^2=\sum_\alpha\alpha!|p_\alpha|^2\).
For rectangular \(Z\), a matrix coefficient
\(p(Z)=w^*W_\nu(Z)u\) has
\begin{equation}\label{eq:selected-fock-matrix}
 \|p\|_{\rm F}^2=H(\nu)\|w\|^2\|u\|^2,\qquad
 H(\nu)=\prod_{(i,j)\in\nu} h_{ij}.
\end{equation}
To see this, realize both Schur modules in tensor powers of
degree \(q=|\nu|\) by Schur--Weyl duality, using the same
unit vector \(s\) in the Specht factor of shape \(\nu\).
The operator \(Z^{\otimes q}\) acts as \(W_\nu(Z)\) on
the Schur factor and as the identity on that Specht factor.
Integrate a coefficient of \(Z^{\otimes q}\) times its
conjugate against independent standard complex Gaussians.

The moment formula pairs positions by all \(\pi\in S_q\),
giving the factor
\(\sum_\pi |\langle s,\pi s\rangle|^2=q!/\dim[\nu]\).
Indeed each pairing identifies the entry in position \(k\)
with the conjugate entry in position \(\pi(k)\), both in
its row and in its column. Indeed the average of the conjugates of the projector onto \(s\)
is scalar on the irreducible Specht module. Its trace is one, so
\[
 \frac{1}{q!}\sum_{\pi\in S_q}\pi P_s\pi^{-1}
     =\frac{1}{\dim[\nu]}I .
\]
Pairing both sides with \(s\) gives the asserted sum. The hook-length
formula gives \(q!/\dim[\nu]=H(\nu)\). This argument also
shows that the Fock normalizer on the square model for the
partition itself is \(H(\nu)\).
\cite{frame-robinson-thrall}

\proofstep{Bounding the homogeneous coefficient.}
For each label rescale to \(z_\tau=r_\tau Z_\tau\) and put
\(g(z)=G((z_\tau/r_\tau)_\tau)\). With \(D=nn'\), take
the torus Fourier projection of \(g\) onto multidegree
\((q_\tau)_\tau\), denoted by \(g_{\mathbf q}\). On the
product of spheres of radius \(\sqrt{q_\tau}\) its
absolute value is at most \(S e^{\sum_\tau q_\tau/2}\).
Use radius zero for a zero degree. Gaussian polar
integration of this homogeneous component yields
\[
 \|g_{\mathbf q}\|_{\rm F}^2
 \le S^2\prod_{\tau:q_\tau>0}
 \frac{e^{q_\tau}(q_\tau+D-1)!}
      {q_\tau^{q_\tau}(D-1)!}.
\]
The radial moment at one label is
\((q+D-1)!/(D-1)!\) because the squared radius of a standard
complex Gaussian vector in \(D\) dimensions has gamma
density with shape \(D\).

For \(q=0\) the factor is one.
For \(q>0\) the factor is at most \(e(q+1)^D\).
Indeed monotonicity of the logarithm gives
\[
 \log q!\le\log q+\int_1^q\log t\,dt
            =\log q+q\log q-q+1 ,
\]
so \(e^q q!/q^q\le eq\), and
\(\binom{q+D-1}{D-1}\le(q+1)^{D-1}\) by recording the first
\(D-1\) entries of a weak composition of \(q\).

\proofstep{Tensor products over the labels.}
For general vectors in the tensor product over labels the
Fock identity \eqref{eq:selected-fock-matrix} tensorizes.
For example, polarizing its proof gives the product of the
input and output Hermitian pairings for two matrix
coefficients. Applying Fock Cauchy--Schwarz to the bilinear
contraction with \(g_{\mathbf q}\) therefore gives the
squared bound \(\|g_{\mathbf q}\|_{\rm F}^2
\prod_\tau H(\nu_\tau)\Omega_\tau^{-q_\tau}\).

To bound \(H(\nu)\), its row product formula gives
\[
 H(\nu)=\frac{\prod_i(\nu_i+n-i)!}
               {\prod_{i<j}(\nu_i-\nu_j+j-i)}
 \le \prod_i \nu_i!\,(q+1)^{n(n-1)/2}.
\]
Here each denominator factor is at least \(j-i\), and
\(\binom{\nu_i+n-i}{n-i}\le(q+1)^{n-i}\).
Concavity of \(\log\) gives
\(k!\le e\sqrt{k}\,\exp(\varphi(k))\) for \(k\ge1\)
by the trapezoid lower bound on its integral from \(1\)
to \(k\). Thus
\(\nu_i!\le(2q+2)\exp(\varphi(\nu_i))\), including
\(\nu_i=0\). Since \(n(n+1)/2\le nn'\), multiplying
these bounds proves \eqref{eq:selected-padded-constant}.
\end{proof}

\paragraph{The finite datum and the two levels.}

The real estimate applies termwise. To use it in the class sum we must specify the finite levels on the two sides of the same entry.

Use the untwisted reference tuple, and fix a word \(u\)
and the two classes of one matrix entry. At the coupled
step also fix their common coset. Put \(g=g_0g_1\),
\(c=T_u g\), and write \(c_B=cc'\). For the integral level exponents \(A,B\) the reference scalars are
\[
 c'=g^{fB-1},\qquad c_A=g^{fA},\qquad fB\ge1 .
\]
Let \(H\) be the semilocal integral Hom module on the two
class lattices at the selected places. In the actual
finite datum \(\Xi\), the factor at the assigned prime is the
indicator of the integral Hom condition there. For a type place \(v\), write \(H_v\) for its integral Hom lattice, \(\mathfrak p_v\) for the maximal ideal and \(f_v\) for the local factor. The two lattice frames give 
\[
 \operatorname{supp} f_v\subseteq H_v,\qquad f_v(x+y)=f_v(x)\quad(y\in\mathfrak p_vH_v).
\]

In ordinary
steps this follows by unfolding the extra columns with
the level character. For a place of residue size \(Q\),
the character has conductor one. If \(u_i,v_i\) are
the ranks including the extra rows, and \(s_i\) counts
extra source columns, every free coordinate in a
contributing matrix has valuation at least \(-s_i\).
At reading depth \(B\) it ranges over at most \(Q^{B+s_i}\)
residue choices. For each free entry the contribution from counting residues and
integrating their cells is bounded by
\[
 Q^{B+s_i}\operatorname{vol}(\varpi^B O)=Q^{B+s_i}Q^{-B}=Q^{s_i},
\]
where \(\varpi\) is a uniformizer and \(O\) has measure one.
The character on the shifted diagonal has argument integral after
multiplication by \(Q\). Hence its value has absolute value one at
every embedding. Multiplying the coordinate estimates and the inverse volume of
the output unipotent gives
\[
 (Q^{B+s_i}Q^{-B})^{u_i v_i}\,Q^{v_i^2}
     =Q^{s_i u_i v_i+v_i^2}\le Q^{2m^2N_m^2}.
\]

Here \(u_i,v_i\le N_m\), \(s_i+1\le m\), and the last factor
on the left is the subgroup index in the padded integral matrix space. In the first coupled entry the
prescribed integral column is inserted into the moved
matrix test. In the last entry the top matrix test is
multiplied by the normalized finite Fourier test on
the last row. Both have residue invariance and bound
one, the latter by the triangle inequality in the
normalized residue average. In both cases the coset
is integral in its decorated frame. The residue sizes
at the type places multiply to \(g^f\). Therefore \(\Xi\)
vanishes outside \(H\), is constant on classes modulo
\(cH\), and, for an absolute \(C\), obeys the uniform bound
\[
      |\rho(\Xi(x))|\le\beta,\qquad \beta=g^{C f m^2 N_m^2}.
\]
\label{cor:selected-kernel-jets-1}

\paragraph{The coefficient field of the two jets.}

A coupled entry involves two jets at once. Their Galois conjugates can be averaged together only after both jets have been placed in one number field.

\begin{lemma}[One field for the two jets in the coupled entry]\label{lem:coupled-jets-field}
At the selected ranks and reference modulus \(M_0=c_Ac_B\),
the chosen coefficient field contains the jets at every permitted
class pair and indicator at each step, including both steps of a
coupled entry. Writing \(K\) for the chosen field and \(a,b\) for any two coupled jet entries, the field of definition of both steps already satisfies
\[
 a,b\in K,\qquad ab\in K.
\]
Their products and class contractions therefore precede the choice of archimedean embedding.
\end{lemma}

\begin{proof}
For the selected ranks \(a,b\), bad-place set \(S_b\), and
positive modulus \(M_0=c_Ac_B\), consider the normalized
indicator indexed by self-dual lattices \(L,L'\),
\(c\mid M_0\), and \(x\in M_0^{-1}H_{S_b}(L,L')\).
Here \(H_{S_b}\) is the semilocal Hom lattice in which
the lift's indicator is defined. For fixed \(L,L',c\),
the indicator depends on \(x\) modulo \(cH_{S_b}\).
The quotient \(M_0^{-1}H_{S_b}/cH_{S_b}\) is finite.

For \(H=H_{S_b}(L,L')\), multiplication by \(M_0\) gives
\[
 M_0^{-1}H/cH\xrightarrow{\ \sim\ }H/(cM_0 H),\qquad
     x+cH\longmapsto M_0 x+cM_0 H.
\]
The local quotient vanishes where \(cM_0\) is a unit.
Use one representative of each of the finitely many
self-dual lattice classes at both ranks. For each pair
and each divisor of \(M_0\), take the normalized Taylor
coefficients of the resulting finite list of indicators
through the total degree of the stage. These coefficients
are algebraic by the CM Taylor comparison on the even-rank
standard lattices and their translates. The hypotheses
are the fixed nonzero period, the compatible ideal character,
and split bad places where the period ideal is a unit.
These hold for the principal curve ideal and the selected
reference sequence.

\proofstep{Changes of representatives.}
Let \(g,h\) identify the source and target lattice representatives with the actual lattices. Composition with these identifications acts on Hom variables by
\[
 x\longmapsto h\circ x\circ g^{-1}.
\] This gives a bijection of the finite semilocal
quotients, since it identifies both \(M_0^{-1}H_{S_b}\)
and \(cH_{S_b}\). Equivariance substitutes the resulting
linear change \(L\) in the labelled Taylor variables.
For a homogeneous term of total degree \(d\), write the
normalized polynomial at the representative as
\(\sum_{|\beta|=d} a_\beta Z^\beta\).

For a coefficient with index \(\alpha\) the transported
reading is
\[
      d_\alpha\sum_{|\beta|=d}a_\beta[Z^\alpha](LZ)^\beta .
\]
Here \(d_\alpha\) is the product of the differential rescaling
in degree \(\alpha\) with the algebraic character ratio.
The latter is a product of determinants to integral powers.
The entries of \(L\) and those determinants use the labelled
images of the auxiliary field. The differential rescalings
use powers of the prescribed generators. Let \(F\) contain the representative coefficients, the entries
of \(L\), the determinant factors and the differential generators.
The multinomial theorem gives
\[
 [Z^\alpha](LZ)^\beta\in \mathbb Z[L_{ij}]\subset F
 \quad\Longrightarrow\quad
 d_\alpha\sum_{|\beta|=d}a_\beta[Z^\alpha](LZ)^\beta\in F.
\] In particular, the coefficients of multidegree
\((|\nu_\tau|)_\tau\), of total degree \(\sum_\tau|\nu_\tau|\),
belong to \(F\).

\proofstep{One finite set for both steps.}
The coefficient-field choice in the selected sequence family
uses this finite generation before the arithmetic row
and column are realized. At fixed predata there are
finitely many words and stages, and finitely many
bad places from the conductors of the character family
and the selected levels. Include the finite generators
for all the permitted subsets of these places.

For the
canonical twisted choice, if \(Q\) is the product of the
row prime norms, its two half moduli are
\(Q^{m+1}\) and \(T_uQ^{m+1}\). Thus the modulus of the stage list is precisely the product
\[
 Q^{m+1}(T_uQ^{m+1})=T_u Q^{2m+2}.
\] The family permits
requiring any such finite collection of algebraic
constants in the embedded coefficient field. Its
construction imposes this requirement jointly with the
column conditions. Thus the chosen field contains the
jets for the actual pair and indicator of each step.
In the coupled entry both consecutive kernels have
their own stage data in this same list. Their products
and the finite class contractions are therefore evaluated
in this field before taking any embedding.
\label{cor:selected-kernel-jets-2}
\end{proof}

\begin{proposition}[Coefficient mass in the torsion presentation]\label{prop:torsion-jet-mass}
At a reference level \(c\) of the selected kernel, use the
torsion presentation of the normalized entry and its finite-data
bound \(\beta\). If \(a_\nu\) are its torsion-presentation coefficients, then for each embedding \(\sigma\)
\[
 \sum_\nu|\sigma(a_\nu)|\le\beta\,c^{2[E:\Q]nn'}.
\]
 This holds also
for a vanishing entry.
\end{proposition}

\begin{proof}
In that field the entry is a jet sum over
\((c_A^{-1}H)/(c_B H)\).
Let \(F_\ell(x)\) be the normalized theta function
attached to the indicator of \(x+\ell H\), using
the class-pair vacuum. For the reference levels, the
odd integer \(c_Ac_B\) is supported on the specified
split half places, where the period ideal is trivial.
To compare the readings in \eqref{eq:unserved-reader}, write
\(A,B\) for the half lattices at \(cc'\), and \(A_1,B_1\)
for those at \(c\). The complementary exponents give
\[
 A\cap B_1\subseteq\Lambda,\qquad
 B_1=(A\cap B_1)+(B\cap B_1),\qquad
 \Lambda\cap A\subseteq A_1 .
\]

In the middle equality the expanding half of each conjugate pair
supplies the local lattice. Put \(D_f=\Lambda\cap A\cap B\)
and \(D_c=\Lambda\cap A_1\cap B_1\). The complementary exponents
also give \(A\cap B\subseteq B_1\), hence \(D_f\subseteq D_c\).
The inclusion \(\Lambda\cap A\subseteq\Lambda\cap A_1\) induces the finite map
\[
 (\Lambda\cap A)/D_f\longrightarrow (\Lambda\cap A_1)/D_c,
 \qquad t+D_f\longmapsto t+D_c .
\]
This map is onto. Passing from \(A_1\) to \(A\) imposes stronger
divisibility only where \(B_1\) expands; there the quotient allows
zero as representative. At the other primes keep the prescribed
residue. The Chinese remainder theorem gives one lift with all
these local residues.

\begin{construction}
\[
 \begin{array}{ccc}
 H/(cc'H)&\longrightarrow&H/cH\\
 x+y&\longmapsto&x \pmod{cH}\\[2pt]
 \{F_{cc'}(x+y):y\in cH/(cc'H)\}&\xrightarrow{\ \mathrm{sum}\ }&F_c(x)
 \end{array}
\]
\caption{Coarsening the Fourier readings in \eqref{eq:selected-kernel-regroup}. The lower arrow sums the fine readings in a single fibre of the upper quotient map.}\label{fig:fine-coarse-fourier}
\end{construction}

\proofstep{Regrouping the finite Fourier series.}
First read the indicator of \(B_1\) in the fine Fourier sum
with no outer translation. The indicator restricts the input
sum to the subgroup \((A\cap B_1)/(A\cap B)\).
For \(y\in B\) its alternating character \(\chi_y\) is already
trivial on \(B\cap B_1\). The sum decomposition of \(B_1\) above
and the self-duality of this lattice therefore give, successively,
\[
 \chi_y|_{A\cap B_1}=1
 \quad\Longleftrightarrow\quad \chi_y|_{B_1}=1
 \quad\Longleftrightarrow\quad y\in B_1 .
\] Orthogonality now leaves exactly the
coarse output representatives in \(B_1/(B_1\cap\Lambda)\).
The two finite groups in the surjection have orders \(h_f,h_c\),
the respective Fourier denominators in \eqref{eq:unserved-reader};
each fibre therefore has size \(h_f/h_c\).
Restrict \(\chi_y\) to
\(B_0=(A\cap B_1)/(A\cap B)\). Then
orthogonality gives
\[
 \frac{1}{|B_0|}\sum_{t\in B_0}\chi_y(t)=\mathbf1_{\{y\in B_1\}} .
\]

At a surviving frequency, counting the fibres of the preceding
surjection gives the coarse normalization \((h_f/h_c)/h_f=1/h_c\).
The semicharacter times the Weyl
translate of the vacuum depends only on the output class,
and the period and volume factors are common.

Expanding the fine reading by its input cosets is thus the
coarse reading at the origin. For a rational lift \(s\) of
the centre, reindex the fine input classes by translation
by \(s\) and act on both sides by \(\alpha(-s)U_{-s}\).
The Weyl composition law and the semicharacter law on \(A\)
give the translated fine indicators with
their same normalization. In semilocal coordinates this is
the partition of the coset modulo \(cH\), which proves
\[
 F_c(x)=\sum_{y\in cH/(cc'H)}F_{cc'}(x+y).
\]
Using the support and constancy of \(\Xi\) gives the
identity of entire sums
\begin{equation}\label{eq:selected-kernel-regroup}
 \sum_{x\in c_A^{-1}H/(cc'H)}\Xi(x) F_{cc'}(x)
        =\sum_{x\in H/cH}\Xi(x) F_c(x).
\end{equation}

In particular the larger denominator bound in the
coefficient field and the refinement factor do not
enter the archimedean count. The projective local
Hom modules have rank \(nn'\) over the local integer
rings, so
\(\#(H/cH)\le c^{[E:\Q]nn'}\). Writing \(q_w\) for the residue size, the count is the partial
product in the norm of \((c)\):
\[
 \prod_{w\ {\rm selected}}q_w^{nn' v_w(c)}
   \le\prod_{w\mid c}q_w^{nn' v_w(c)}
   = |N_{E/\Q}(c)|^{nn'} .
\] The missing factors, when present,
are positive integral powers in the full norm of \(c\).

\proofstep{Bounding the number and sizes of the coefficients.}
At level \(c\), write \(s\) for the rational lift of \(x\)
in the half lattice and use \(S_A,S_B\) as in the
finite Fourier reading. The expansion of \(F_c(x)\)
uses translates at \(y-s\), \(y\in S_B\). With \(v\) the
volume in the indicator transform, the two factors
\(v^{-1/2}\) cancel the \(v\) in its Fourier measure.
The phase multiplying the average at \(y\) is
\(\alpha(s)\alpha(y)e^{\pi i E_H(s,y)}\), since the
indicator translation is \(U_{-s}\). Write \(B_{\rm lat}\)
for the contracting half lattice. The average is
\[
 \frac1{|S_A|}\sum_{t\in S_A}{\bf1}_{B_{\rm lat}}(t)
                                e^{2\pi i E_H(t,y)}.
\]
All its phases, including the prefactor at \(y\),
are roots of unity by rationality of \(E_H\) and the
quadratic function defining \(\alpha\) on the Hom space.
Thus each coefficient is algebraic and of absolute
value at most one at every embedding.

Multiplication
by \(c\) sends the period lattice into the contracting
half lattice, so \(|S_B|\le c^{[E:\Q]nn'}\). The rational
lifts \(y-s\) give torsion points on the period torus.
Together with \eqref{eq:selected-kernel-regroup} this
bounds the sum of conjugate coefficient moduli in the
torsion presentation by
\begin{equation}\label{eq:selected-kernel-coeff-l1}
             \beta\,c^{2[E:\Q]nn'} .
\end{equation}
The ideal character used to define the moved vacuum
has been divided out in this presentation. The
remaining ratio of frame determinant characters has
modulus one at each embedding for the reference level.
Indeed all frame changes here are at the type primes,
where their uniformizer factors cancel against the tied
normalized local character. The surviving residue
characters have finite image. A zero entry satisfies the
bound directly.
\label{cor:selected-kernel-jets-3}
\end{proof}

\paragraph{The Gaussian bound after conjugation.}
Fix an inducing embedding \(\rho\) of the defining field.
The normalized theta vacuum in the two class lattices
is the moved vector of Section~\ref{subsec:theta-normalization}.
Here the two ranks are even, the period lattice is self-dual,
and the chosen nonzero two-torsion point survives multiplication by both
\((1+\sqrt{-d})/2\) and \((1-\sqrt{-d})/2\). The
coefficientwise conjugation of the one-variable normalized
theta and its torsion translates was described in
Section~\ref{subsec:selected-period-means}. It applies both
to the finite sum transporting the vacuum and to each
further rational translate in the indicator expansion.
Equivalently, the torsion transport acts on theta operators
with their Weil pairing.

If the phase changes by a root of unity \(\zeta\), the pointwise absolute bound is preserved, since
\[
 |\zeta a|=|\zeta|\,|a|=|a|.
\]
The induced CM type permutes the labels under embeddings
fixing the quadratic field. For an embedding exchanging
its two embeddings, also apply complex conjugation to the
coordinates. This gives entire conjugate functions in the
same labelled Euclidean metric, with the conjugate period
area \(A_\rho\), and preserves the squared average of the
vacuum. If the normalized sum is \(\sum_a c_aV_a\), its distinct
translate characters are orthogonal, also after conjugation. Thus
\[
 \left\|\sum_a\zeta_a c_aV_a\right\|^2
   =\sum_a |c_a|^2\,\|V_a\|^2\qquad(|\zeta_a|=1).
\] In particular the conjugate norm
squared equals \(N_\rho^{D_j}\) for
\(D_j=f nn'\), as computed in
Section~\ref{subsec:selected-period-means}.
\label{cor:selected-kernel-jets-4}

\begin{lemma}[Gaussian prefactor at a conjugate embedding]\label{lem:conjugate-jet-prefactor}
For the selected even ranks and a theta vacuum with self-dual
period lattice, assume its marked two-torsion point has nonzero
images under both \((1+\sqrt{-d})/2\) and \((1-\sqrt{-d})/2\).
The single-kernel Gaussian bound at an inducing embedding \(\rho\)
has a prefactor \(S\) satisfying
\[
 S^2\le N_\rho^{D_j}(8+16D_jr_E)^{D_j} e^{2(C+4)U_u}.
\]
The bound applies to either kernel at the coupled step.
\end{lemma}

\begin{proof}
Here is the pointwise bound with the covolume normalization.
Let \(r_E\) be the root discriminant. Write the period
lattice as the embedded module
\(\delta\,\mathfrak c I\operatorname{Hom}(L,L')\),
where \(\mathfrak c\) is the integral curve ideal of the
quadratic field, \(\delta\) is the period scale, and
\(I\bar I=\mathfrak D_{E/K}^{-1}\).
For a nonzero \(x\in\mathfrak c I\operatorname{Hom}(L,L')\)
the relative matrix trace \(v=\operatorname{tr}(x^*x)\)
is totally positive in the real subfield. Local frames
of the two self-dual modules show that
\(x^*x L\subset \mathfrak c\bar{\mathfrak c}I\bar I L\),
so this trace lies in
\(\mathfrak c\bar{\mathfrak c}I\bar I\).

Taking absolute ideal norms and using arithmetic-geometric
mean over the \(f\) real embeddings gives
\[
 \sum_\tau \tau(v)
    \ge f\,\mathrm N_{E/\Q}
                 (\mathfrak c\bar{\mathfrak c}I\bar I)^{1/(2f)}
    = f\,\mathrm N_{K/\Q}(\mathfrak c)\,\sqrt d/r_E .
\]
The area of the curve lattice is
\(|\delta|^2\mathrm N(\mathfrak c)\sqrt d/2\).
This gives minimum squared length at least
\(2f A_\rho/r_E\) also for the conjugate period lattice.
Its covolume is \(A_\rho^{D_j}\), by self-duality for
the imaginary part of the Hermitian form
\(H_\rho(z,z)=\|z\|^2/A_\rho\).

\proofstep{The submean estimate on the period torus.}
Let \(V\) be this conjugate theta vector or a Heisenberg
translate of it. Its weighted squared absolute value
\(|V(z)|^2 e^{-\pi H_\rho(z,z)}\) is periodic.
Translation by \(U_w\) translates this density and preserves
its mean. The ball of radius
\(R=(A_\rho/(2r_E+1))^{1/2}\) embeds in the period torus. Indeed for the shortest nonzero
period length \(\lambda_1\),
\[
 (2R)^2=\frac{4A_\rho}{2r_E+1}
    <\frac{2A_\rho}{r_E}\le\lambda_1^2 .
\]
Two points of the ball therefore cannot differ by a nonzero period.

Apply the submean inequality
at zero to \(U_z V\) on that ball, then insert the Gaussian
weight on the ball. Since its volume is
\(\pi^{D_j}R^{2D_j}/D_j!\), the squared bound at \(z\)
is
\[
 \frac{D_j!}{\pi^{D_j}}(2r_E+1)^{D_j}
    e^{\pi/(2r_E+1)}
    N_\rho^{D_j}\,e^{\pi\|z\|^2/A_\rho}
 \le 3(2D_jr_E+1)^{D_j}
    N_\rho^{D_j}\,e^{\pi\|z\|^2/A_\rho}.
\]
The last inequality uses \(D_j\ge1\) and \(r_E\ge1\).
This bound covers all classes and all the translates
in \eqref{eq:selected-kernel-coeff-l1}.

\proofstep{Summing the torsion translates.}
Apply the triangle inequality to their conjugate entire
sum. It satisfies the Gaussian hypothesis above with
\[
 S=\beta\,c^{2[E:\Q]nn'}\sqrt{3(2D_jr_E+1)^{D_j}}\,
                                  N_\rho^{D_j/2},
 \qquad r_\tau=\sqrt{\pi/A_\rho}.
\]
The Taylor differential multiplier in label \(\tau\)
is the conjugate of the generator of the coordinate axis
\(d_\tau\). The normalized metric is consequently
\(\Omega_\tau=A_\rho/(\pi|\rho(d_\tau)|^2)\),
so \(|r_\tau\rho(d_\tau)|^2=\Omega_\tau^{-1}\).
The formal jet in the coefficient field agrees with
this conjugate analytic jet.

Each relevant monomial has
the prescribed total degree at each label, and comparison
before applying \(\rho\) is by the multivariable Taylor
coefficients, multiplied by their multi-index factorials
and the labelled differential powers. The complex
coefficient form is the Gaussian pairing normalized on
the highest polynomial in each label as in
Lemma~\ref{lem:padded-jet-estimate}. Thus the
same operator bound holds on the embedded coefficient
spaces. Finally \([E:\Q]=2f\) and
\(D_j\le fm^2N_m^2\). Using \(3\le 8^{D_j}\) gives
\[
 S^2\le N_\rho^{D_j}(8+16D_jr_E)^{D_j}
                       e^{2(C+4)U_u}.
\]
This is the single-entry prefactor used in the conjugate
mean, including both entries of the last coupled step.
\label{cor:selected-kernel-jets-5}
\end{proof}
